\section*{第 \ref{ch:b2:aromatic-substitution} 章——芳香性与亲电芳香取代}

\begin{solution}{exo:b2:aromatic-substitution:1}
\ce{HNO3 + 2H2SO4 -> NO2+ + H3O+ + 2HSO4-}：硝酸被质子化，然后失去水。
\end{solution}

\begin{solution}{exo:b2:aromatic-substitution:2}
甲苯生成 2-溴甲苯和 4-溴甲苯：甲基定位至邻位和
对位。硝基苯生成 3-溴硝基苯，且反应更慢：硝基定位至
间位。
\end{solution}

\begin{solution}{exo:b2:aromatic-substitution:3}
\ce{-CH3}：活化，邻/对位定位。\ce{-OH}：强活化，邻/对位定位。\ce{-Cl}：钝化，邻/对位定位。
\ce{-NO2}：强钝化，间位定位。\ce{-COCH3}：钝化，间位定位。\ce{-NHCOCH3}：
中等活化，邻/对位定位。
\end{solution}

\begin{solution}{exo:b2:aromatic-substitution:4}
苯乙酮（1-苯基乙-1-酮），\ce{C6H5COCH3}，经酰基正离子 \ce{CH3CO+} 生成。
\end{solution}

\begin{solution}{exo:b2:aromatic-substitution:5}
伯碳正离子（或其配合物）经氢负离子迁移，
重排为仲碳正离子（异丙基正离子），再使苯环烷基化。丙苯：先以丙酰氯进行傅–克酰基化（无重排），
再将酮的 \ce{C=O} 还原为 \ce{CH2}。
\end{solution}

\begin{solution}{exo:b2:aromatic-substitution:6}
\ce{-OH} 为强共轭给电子基团：苯环活化强，分子溴
即可作为足够的亲电试剂，邻、对位三个位置全部被取代。
\end{solution}

\begin{solution}{exo:b2:aromatic-substitution:7}
3-溴硝基苯：先硝化，再溴化（硝基定位至间位）。4-溴硝基苯：
先溴化，再硝化（溴定位至邻、对位），然后将对位异构体
与邻位异构体分离。
\end{solution}

\begin{solution}{exo:b2:aromatic-substitution:8}
混酸会氧化电子云密度很高的苯胺，并使其氮原子质子化：\ce{-NH3+}
是间位定位的钝化基团。在乙酰苯胺中，氮上的孤对电子与羰基共享，
碱性和活化作用均减弱，既不会被质子化，也不会被氧化，
仍定位至邻、对位，且因空间位阻，以对位为主。
\end{solution}

\begin{solution}{exo:b2:aromatic-substitution:9}
先磺化甲苯（以对位为主，体积较大的基团避开邻位）；再溴化：溴进入甲基的邻位
（对位已被占据，两个基团均有利于该位置）；
用稀酸加热除去 \ce{SO3H} 基团。
\end{solution}

\begin{solution}{exo:b2:aromatic-substitution:10}
两个基团均定位至邻、对位，其对位各被对方占据。甲氧基
活化作用较强，起主导作用：硝化发生在 \ce{-OCH3} 的邻位（2 位），
生成 4-甲基-2-硝基苯甲醚。
\end{solution}

\begin{solution}{exo:b2:aromatic-substitution:11}
每个硝基都强烈钝化苯环：第二次硝化比第一次需要更强的条件，
第三次（苯环上已有两个硝基，仅余弱活化的甲基）所需条件
则更为苛刻。
\end{solution}

\begin{solution}{exo:b2:aromatic-substitution:12}
4-硝基苯甲酸：先硝化甲苯，分出对位异构体，再用热高锰酸钾将甲基
氧化为 \ce{-COOH}。3-硝基苯甲酸：先将甲苯氧化为苯甲酸，再硝化
（\ce{-COOH} 定位至间位）。
\end{solution}

\begin{solution}{pb:b2:aromatic-substitution:1}
\textbf{1.} \ce{HNO3 + 2H2SO4 -> NO2+ + H3O+ + 2HSO4-}.
\textbf{2.} \ce{O=N+=O}：直线形，与 \ce{CO2} 等电子。
\textbf{3.} 氮与苯环上位于 \ce{CH3} 对位的碳成键，该碳变为四面体形
（连有 H 和 \ce{NO2}）。
\textbf{4.} 正电荷分布在被进攻碳的两个邻位碳，以及与之处于对位、
连有甲基的碳上。
\textbf{5.} 第一步，生成惠兰德中间体。
\textbf{6.} 介质中的碱：\ce{HSO4-} 或水。
\textbf{7.} 邻位两个、间位两个、对位一个。
\textbf{8.} 2 : 2 : 1，即邻位 40~\%、间位 40~\%、对位 20~\%。
\textbf{9.} 间位进攻时，正电荷从不位于连有甲基的碳上，
没有任何共振结构得到甲基的帮助。
\textbf{10.} 有一个共振结构将正电荷置于连有 \ce{CH3} 的碳上：叔碳正离子，
由甲基的诱导效应和超共轭作用而稳定。
\textbf{11.} 邻、对位有利（96~\%，而随机情况下为 60~\%）；间位几乎不生成。
\textbf{12.} 按位置计：邻位 $58/2 = 29~\%$，对位 38~\%：对位是活性更高的位置，
邻位因甲基而略受阻碍。
\textbf{13.} 更快：甲基活化苯环。
\textbf{14.} 4-硝基甲苯（熔点为 \qty{52}{\degreeCelsius}）。
\textbf{15.} 冷却粗产物混合物：对位异构体结晶析出并过滤除去；
液体中保留邻位和间位异构体。
\textbf{16.} 对称分子在晶体中堆积更紧密：单位体积内相互作用更多，
因而熔点更高。
\textbf{17.} 可用减压分馏，或用色谱法分离；也可
进一步冷却（间位异构体熔点为 \qty{16}{\degreeCelsius}）。
\textbf{18.} 异构体的沸点相近（分子式相同，极性相近）。
\textbf{19.} \ce{C7H8}: \qty{92.14}{g/mol}; \ce{C7H7NO2}: \qty{137.14}{g/mol}.
\textbf{20.} 甲苯的物质的量为 $100/92.14 = \qty{1.085}{mol}$；硝基甲苯
的物质的量为 $0.90 \times 1.085 = \qty{0.977}{mol}$。
\textbf{21.} 总计 $0.977 \times 137.14 = \qty{134.0}{g}$：邻位 \qty{77.7}{g}，间位
\qty{5.4}{g}，对位 \qty{50.9}{g}。
\textbf{22.} 再进行一次硝化（生成二硝基甲苯），以及甲基的氧化。
\textbf{23.} 用热碱性高锰酸钾将甲基氧化为 \ce{-COOH}，然后酸化。
\textbf{24.} 本次操作可得 $\boldsymbol{\approx \qty{51}{g}}$ 的 4-硝基甲苯。
\end{solution}
